\documentclass[11pt,a4paper]{article}
\usepackage[margin=18mm,top=17mm,bottom=19mm]{geometry}
\usepackage{amsmath,amssymb,graphicx,array,fancyhdr,lastpage,textcomp}
\usepackage{fontspec}
\setmainfont{texgyreheros-regular.otf}[BoldFont=texgyreheros-bold.otf,ItalicFont=texgyreheros-italic.otf,BoldItalicFont=texgyreheros-bolditalic.otf]
\setlength{\parindent}{0pt}
\setlength{\parskip}{3pt}
\pagestyle{fancy}\fancyhf{}
\renewcommand{\headrulewidth}{0pt}\renewcommand{\footrulewidth}{0.3pt}
\fancyfoot[L]{\footnotesize by paper.shrit.in\quad |\quad normal}
\fancyfoot[R]{\footnotesize Page \thepage\ of \pageref{LastPage}}
\newcommand{\question}[3]{\par\vspace{8pt}\noindent\begin{minipage}[t]{0.075\linewidth}\textbf{#1)}\end{minipage}\begin{minipage}[t]{0.855\linewidth}#3\end{minipage}\hfill\begin{minipage}[t]{0.055\linewidth}\raggedleft(#2)\end{minipage}\par}
\begin{document}
{\LARGE\bfseries Question Paper}\hfill {\small Registration No.: \rule{33mm}{0.3pt}}\par\vspace{8pt}
\begin{center}{\large\bfseries MANIPAL FORMAT | PRACTICE PAPER}\\[6pt]{\bfseries THIRD SEMESTER B.TECH. | MIDSEM PRACTICE}\\[4pt]{\bfseries DATA ANALYTICS [CSS 2103]}\\[4pt]{\small COMPUTER SCIENCE AND ENGINEERING}\\[3pt]{\small SET A: ENSEMBLE OF PAST-PAPER QUESTIONS}\end{center}
\textbf{Marks: 30}\hfill\textbf{Suggested duration: 90 minutes}\par\vspace{3pt}\hrule\vspace{5pt}
{\small\textbf{Answer all questions.} Questions 1--5 are MCQs worth 1 mark each. Select one correct option per MCQ. The remaining questions are theory questions worth 25 marks in total; show working and draw labelled diagrams where appropriate. Modules 1-3. Independent practice paper; not an official university examination.}\par
{\small Non-programmable calculator permitted. Use the critical values supplied.}\par
\medskip\textbf{Section A: Multiple-choice questions (5 marks)}\par
\question{1}{1}{A manufacturer claims a population mean accuracy of 0.95. Which null hypothesis tests this claim?\par (A) $\mu=0.95$\par (B) $\bar x=0.95$\par (C) $\mu\ne0.95$\par (D) $\mu<0.95$}
\question{2}{1}{A population has known standard deviation 0.02. For an independent sample of size 36, what is the standard error of the sample mean?\par (A) 0.12\par (B) 0.02\par (C) $0.02/6$\par (D) $0.02/36$}
\question{3}{1}{With $\bar x=0.947$, $\mu_0=0.95$, $\sigma=0.02$ and $n=36$, what is the $z$ statistic?\par (A) 0.90\par (B) -0.90\par (C) -9.00\par (D) 9.00}
\question{4}{1}{Observations range from 180 to 480. What is the width of each of three equal-width bins?\par (A) 60\par (B) 150\par (C) 300\par (D) 100}
\question{5}{1}{Twelve observations are divided into three equal-frequency bins. How many observations are in each bin?\par (A) 3\par (B) 4\par (C) 6\par (D) 12}
\newpage\textbf{Section B: Theory questions (25 marks)}\par
\question{6}{5}{(i) Find the standard deviation and skewness for: 78, 91, 68, 54, 82, 47, 60, 78, 97, 58, 62, 50, 50, 63, 92, 75, 79, 63, 42, 61.\par (ii) Give an example of the combining-variables concept in dataset preparation.\par (iii) A CMO compares three new cola flavours using ratings from 12 customers on a 1-9 scale. Conduct the appropriate analysis. Ratings: Flavour A: 4, 2, 4, 2; Flavour B: 6, 4, 5, 5; Flavour C: 8, 6, 8, 6. Use population standard deviation and Pearson's second skewness coefficient. For (iii), assume independent normal samples with equal variances; use a one-way ANOVA at 5\%, with critical F(2,9) = 4.256.}
\question{7}{5}{(i) Find the elements of a box plot for: 78, 91, 68, 54, 82, 47, 60, 78, 97, 58, 62, 50, 50, 63, 92, 75, 79, 63, 42, 61.\par (ii) Why is it sometimes necessary to create subsets from a large dataset before analysis?\par (iii) Measure the strength of the relationship between the rankings assigned by two interviewers to candidates A-L. Interviewer 1: 1, 2, 3, 4, 5, 6, 7, 8, 9, 10, 11, 12. Interviewer 2: 1, 2, 4, 3, 6, 5, 8, 7, 10, 9, 12, 11. For quartiles, use medians of the two halves. Use Spearman rank correlation in (iii).}
\question{8}{3}{A hospital laboratory claims that a new COVID testing kit correctly identifies infection 95\% of the time. A researcher believes the true accuracy is less than 95\%. A sample of 50 test results has mean accuracy 0.93; the known population standard deviation is 0.04. Test the claim at the 5\% significance level. Critical z-value: -1.645.}
\question{9}{5}{Test the effectiveness of vaccination in preventing smallpox at the 5\% significance level.\par Vaccinated: attacked 31, not attacked 469, total 500.\par Not vaccinated: attacked 185, not attacked 1315, total 1500.\par Overall: attacked 216, not attacked 1784, total 2000. Use Pearson chi-square without continuity correction; critical value for 1 degree of freedom is 3.841. State association, not a causal conclusion.}
\question{10}{5}{Per-acre wheat production (tons) is recorded for three varieties and four fertilizers, with each combination replicated twice. Test whether variety differences and fertilizer differences are significant at the 5\% level. Each pair below gives the two replicates, in fertilizer order F1, F2, F3, F4.\par Variety 1: (8,9), (9,8), (6,7), (10,11).\par Variety 2: (7,6), (8,9), (5,4), (9,10).\par Variety 3: (10,9), (11,12), (8,7), (12,13). Assume independent normal errors with common variance. Include the interaction term. Use critical F values: variety (2,12): 3.885; fertilizer (3,12): 3.490.}
\question{11}{2}{A data scientist is analysing the following cholesterol level (mg/dL) readings of 12 patients: [145, 155, 160, 162, 165, 172, 180, 190, 200, 220, 240, 260]. Apply equal-width and equal-frequency binning with 3 bins. Using histogram analysis, explain which binning approach (width-based or frequency-based) better represents the data distribution and why.}
\vfill\begin{center}\small End of question paper\end{center}\end{document}